\chapter{Programmable Hardware II: Trading Designs}\label{ch:nw:programmable-hardware-ii-trading-designs}

The design of this chapter decides two clock cycles after the last byte of a price reaches it: \qty{12.8}{\nano\second} at 156.25~MHz,
for the first order and for the ten-thousandth, whether the market is quiet or bursting. It has no cache to miss and no thread to wake,
and the order it sends leaves as a template with three fields patched in. It also cannot do much: it compares one price with one
threshold, for one instrument, and everything it knows about the world was written into its registers by software a moment before.
That combination (very fast, very narrow, steered by software) is what firms mean by a \omterm{def:nw:programmable-hardware-ii-trading-designs:hybrid}{hybrid trading system}.

This chapter builds a \omterm{def:nw:programmable-hardware-ii-trading-designs:trigger}{hardware trigger} on the exchange simulator's feed, with its risk checks and its \omterm{def:nw:programmable-hardware-ii-trading-designs:template}{order template}, in SystemVerilog
on the toolkit of chapter~6, and checks it four ways: in Verilator, in Icarus, against a C++20 cycle model and against a Python model,
order for order and cycle for cycle. It then puts the trigger in the firm's wire-to-wire path and asks what it buys.

\section{Feed decoding in hardware}

A software handler (One Quant Book~13, chapter~18) receives a whole datagram, then walks its messages. A hardware decoder walks the
messages as the bytes arrive, eight a cycle, keeping only the state it needs between cycles: where it is in the packet's header, the
length of the current message, how far into it it is, and the fields it has captured so far. In the simulator's feed, as in the
exchange feed format it follows, an add-order message is 36 bytes with its price in the last four; the decoder therefore knows everything
about an order the cycle its last byte arrives, while the rest of the packet (other messages, the frame's check sequence) may still be to
come.

\omcode{code/firm/hwtrade/hdl/hwt_trigger.sv}{47}{75}{The decoder's combinational logic: a byte-serial state machine unrolled over the
beat's eight lanes. It captures the type, instrument, side, size and price, and at the message's last byte decides whether to
fire.}\label{lst:nw:programmable-hardware-ii-trading-designs:parse}

\begin{definition}[Hardware trigger, cut-through decision]\label{def:nw:programmable-hardware-ii-trading-designs:trigger}
A \emph{hardware trigger}\index{hardware trigger} is logic on a network card that watches a market-data stream for a condition set in
advance by software (a price crossing a threshold, a trade on a given instrument) and, when it occurs, starts an order without the
host's involvement. A \emph{cut-through decision}\index{cut-through decision} is a decision taken as soon as the bytes it depends on
have arrived, before the frame that carries them has ended or been checked.
\end{definition}

The unrolled state machine of \cref{lst:nw:programmable-hardware-ii-trading-designs:parse} is the simplest correct design, and it is
slow in hardware: eight byte steps chained in one cycle make a long \omterm{def:nw:programmable-hardware-i-architecture-and-toolchain:timing}{critical path}. Faster decoders exploit the format: messages start at
known offsets relative to their length prefixes, so dedicated logic can look for a message start in each lane at once, and fixed-layout
messages can be matched by comparing whole words. That is the work of chapter~6's \omterm{def:nw:programmable-hardware-i-architecture-and-toolchain:timing}{timing closure}, and it is why production decoders are
specific to one venue's format.

\section{Book building in hardware}

A trigger on single messages needs no book. A strategy that trades on the best bid or offer needs one, and a book in hardware is a
different data structure from the ones of One Quant Book~13, chapter~19. Orders by reference need a hash table in \omterm{def:nw:programmable-hardware-i-architecture-and-toolchain:bram}{block RAM} or external
memory; price levels need an array indexed by the price's offset from a reference, with the best level found by a priority encoder over
occupied levels; and every update must finish within the time between messages, or the design must queue them. Firms that build books
in hardware keep only a few levels near the touch, for the few instruments that need them, and leave the rest to software.

\section{Risk checks in cycles}

A firm's hardware sends orders on behalf of a broker-dealer that must keep pre-trade controls ``under the direct and exclusive
control'' of the broker (One Quant Book~11, chapter~27), controls that reject orders exceeding ``appropriate price or size parameters,
on an order-by-order basis or over a short period of time''. In hardware those checks are comparisons and counters, and they cost a
pipeline stage.

\begin{definition}[Register map]\label{def:nw:programmable-hardware-ii-trading-designs:regmap}
A \emph{register map}\index{register map} is the set of addresses at which a hardware design exposes its configuration and state to
software (thresholds, limits, enable and kill bits, counters), read and written over the host's bus, and the document that defines
them.
\end{definition}

\omcode{code/firm/hwtrade/hdl/hwt_trigger.sv}{108}{135}{The risk stage and the order template: a size limit, a kill bit and a token
bucket, then the order's id, size and price patched in.}\label{lst:nw:programmable-hardware-ii-trading-designs:risk}

\begin{proposition}[A token bucket's burst]\label{prop:nw:programmable-hardware-ii-trading-designs:bucket}
A token bucket of capacity $B$ refilled with one token every $R$ cycles lets out at most $B + \lfloor t/R\rfloor$ orders in any interval
of $t$ cycles. If software needs $t$ cycles to notice a runaway and set the kill bit, the hardware can send at most that many orders
before it stops, and none after.
\end{proposition}

\begin{proof}
The bucket holds at most $B$ tokens at the start of the interval and gains at most one every $R$ cycles; each order consumes one. The
kill bit masks the output stage in the cycle it is set.
\end{proof}

The proposition is the requirement document of a hardware risk layer. At 156.25~MHz, with the chapter's bucket of 4 tokens refilled
every 64 cycles, a software monitor that reacts in \qty{20}{\micro\second} lets out at most 52 orders; one that reacts in
\qty{100}{\micro\second}, 248. The limits that matter in hardware are those that bind before software can act.

\section{Order generation from templates}

\begin{definition}[Order template]\label{def:nw:programmable-hardware-ii-trading-designs:template}
An \emph{order template}\index{order template} is an order message prepared in advance, complete except for the fields that depend on
the trigger (client order identifier, size, price), which the hardware patches in as it sends.
\end{definition}

\begin{definition}[TCP offload engine]\label{def:nw:programmable-hardware-ii-trading-designs:toe}
A \emph{TCP offload engine}\index{TCP offload engine} implements the TCP protocol (sequence numbers, acknowledgements, retransmission)
in a network card's hardware, so that the card can send and receive on a TCP connection without the host.
\end{definition}

Order entry runs over TCP (chapter~1), so a card that sends orders by itself must own the connection's state: its next sequence
number, its acknowledgements, its retransmissions. A \omterm{def:nw:programmable-hardware-ii-trading-designs:toe}{TCP offload engine} does that; the \omterm{def:nw:programmable-hardware-ii-trading-designs:template}{order template} then carries the TCP and IP
headers too, with sequence number and checksums patched at send time. The software that owns the session (logon, heartbeats, the
venue's session layer) shares the connection with the hardware through the same engine.

\begin{dated}{2026-09}{A published hardware tick-to-trade figure}\label{dat:nw:programmable-hardware-ii-trading-designs:stac}
In a STAC-T0 benchmark reported in June 2024, an FPGA system built on an AMD Alveo UL3524 card with a TCP and UDP core (an Exegy and AMD
solution) was measured at a minimum of \qty{13.9}{\nano\second} for 507-byte frames and \qty{14.1}{\nano\second} for 68-byte frames, from
``the last bit of inbound data needed to make a trading decision to the first bit of the simulated outbound order''. The definition
excludes the time for the inbound bits to arrive, which the chapter's design, and any design, cannot avoid.
\end{dated}

\section{Hybrid designs: what stays in software and why}

\begin{definition}[Hybrid trading system]\label{def:nw:programmable-hardware-ii-trading-designs:hybrid}
A \emph{hybrid trading system}\index{hybrid trading system} splits a strategy between software, which computes its model, positions
and parameters at its own pace, and hardware, which applies precomputed decisions to the market data at \omterm{def:nw:networking-for-trading:ser}{line rate}; software writes the
hardware's registers and reads its counters and copies of what it sent.
\end{definition}

The split follows from what each side is good at. Hardware is fast and fixed: every feature costs logic, every change costs a build of
hours and a verification of days, and it cannot hold much state. Software is slow and flexible: it can model, learn, net positions
across venues and change within minutes. So software decides \emph{what} would be worth doing (``buy up to 700 at 100.02 or better if
anyone offers'') and hardware decides \emph{when} (the cycle the offer arrives). Everything that needs history, many instruments or
judgement stays in software; the one comparison that must happen first moves to the card.

\begin{omfigure}
\begin{tikzpicture}[font=\footnotesize,b/.style={draw,rounded corners=2pt,minimum height=0.7cm,align=center,inner sep=3pt},>=Stealth]
\node[b,fill=gray!10] (feed) at (0,0) {feed\\ (UDP)};
\node[b,fill=omDef!18] (dec) at (2.5,0) {decode\\ (stage 1)};
\node[b,fill=omThm!15] (risk) at (5.4,0) {risk: size,\\ bucket, kill\\ (stage 2)};
\node[b,fill=omMeth!20] (tpl) at (8.3,0) {order\\ template\\ (stage 3)};
\node[b,fill=gray!10] (out) at (11.0,0) {orders\\ (TCP offload)};
\node[b,fill=omDef!8,minimum width=5.6cm] (sw) at (4.0,-2.4) {software on the host: model, positions, limits};
\node[b,fill=gray!8] (regs) at (4.0,-1.25) {register map: thresholds, limits, kill, counters};
\draw[->,thick] (feed) -- (dec);
\draw[->,thick] (dec) -- node[above,font=\scriptsize]{trigger} (risk);
\draw[->,thick] (risk) -- (tpl);
\draw[->,thick] (tpl) -- (out);
\draw[->,omThm] (sw) -- (regs);
\draw[->,omThm] (regs.north -| risk.south) -- (risk.south);
\draw[->,omThm] (regs.north -| dec.south) -- (dec.south);
\draw[->,gray,dashed] (tpl.south) |- node[pos=0.25,right,font=\scriptsize]{copies of orders} (sw.east);
\end{tikzpicture}

\omcaption{The chapter's hybrid design. Three pipeline stages on the card decode the feed, check each candidate order and patch it into
a template; software on the host writes the thresholds, limits and kill bit through the \omterm{def:nw:programmable-hardware-ii-trading-designs:regmap}{register map} and receives copies of everything
the card sends.}\label{fig:nw:programmable-hardware-ii-trading-designs:hybrid}
\end{omfigure}

\begin{omfigure}
\begin{tikzpicture}[font=\scriptsize,c/.style={draw,minimum height=0.55cm,minimum width=1.05cm,align=center,inner sep=1pt}]
\foreach \i in {0,...,9} {\node at (1.1*\i,1.25) {\i};}
\node[anchor=east] at (-0.7,1.25) {cycle};
\foreach \i in {0,...,7} {\node[c,fill=omDef!15] at (1.1*\i,0.55) {beat \i};}
\node[anchor=east] at (-0.7,0.55) {input};
\node[c,fill=omDef!35] at (1.1*7,-0.15) {price};
\node[anchor=east] at (-0.7,-0.15) {decode};
\node[c,fill=omThm!25] at (1.1*8,-0.85) {risk};
\node[anchor=east] at (-0.7,-0.85) {risk};
\node[c,fill=omMeth!30] at (1.1*9,-1.55) {order};
\node[anchor=east] at (-0.7,-1.55) {template};
\end{tikzpicture}

\omcaption{One packet with one add-order message (58 bytes: 20 of header, 2 of length, 36 of message) through the design, one beat of 8
bytes a cycle. The price's last byte arrives in beat~7; the trigger is registered at the end of that cycle, the risk decision one cycle
later, and the order in the cycle after: two cycles after the price, ten after the packet's first byte.}\label{fig:nw:programmable-hardware-ii-trading-designs:cycles}
\end{omfigure}

\section{Tutorial: a trigger on the simulator's feed, four ways}\label{tut:nw:programmable-hardware-ii-trading-designs:tut}

\begin{tutorial}
\textbf{Goal.} Run the trigger on two seconds of the exchange simulator's feed in Verilator, in Icarus, in a C++20 cycle model and in a
Python model, check that all four send the same orders on the same cycles, and put the design in the wire-to-wire path.
\textbf{End state:} \cref{tab:nw:programmable-hardware-ii-trading-designs:runs}, \cref{fig:nw:programmable-hardware-ii-trading-designs:race}
and the green tests of \texttt{firm.hwtrade}.
\begin{tutsteps}
\item \textbf{The feed.} \texttt{make\_hwtrade\_fixture.py} runs One Quant Book~10's simulator for two seconds from the open (992 packets)
  and stores them; \texttt{beats\_of} cuts each into beats with two idle cycles between packets (the day compressed: the bucket counts
  cycles, not seconds).
\item \textbf{Four models.} \texttt{build\_verilator}, \texttt{build\_icarus}, \texttt{CycleModel} and \texttt{cpp/firm\_hwtrade.hpp}
  (\cref{lst:nw:programmable-hardware-ii-trading-designs:model}) run the same settings; the tests compare their outputs line for line.
\item \textbf{The reference.} \texttt{triggers()} parses the same packets message by message in Python: every qualifying message must be
  either an order or a reject.
\item \textbf{The path.} \texttt{nw\_hwtrade.paths()} builds the hardware variant of \texttt{firm.wirepath} (the cage of chapter~2, the
  card's transceivers, three cycles) next to chapter~5's software path, and \texttt{win\_probability} races each against a competitor;
  \texttt{python fig\_hwtrade.py} writes the chart's data.
\end{tutsteps}
\textbf{What to change next.} Set the bucket to one token refilled every 10\,000 cycles and count the orders; move the kill bit to the
cycle before the tenth order and check that exactly nine leave.
\end{tutorial}

\omcode{code/firm/hwtrade/cpp/firm_hwtrade.hpp}{20}{44}{The C++20 cycle model's clock edge: stage 3, then stage 2, then stage 1, each
reading the registers the previous edge wrote, as the hardware does.}\label{lst:nw:programmable-hardware-ii-trading-designs:model}

\begin{table}[ht]
\centering\small
\begin{tabular}{rrrrr}
\toprule
Threshold (price $\times 10^4$) & Maximum size & Kill set at cycle & Orders & Rejects \\
\midrule
1\,000\,200 & 10\,000 & never & 108 & 5 \\
1\,000\,300 & 10\,000 & never & 133 & 27 \\
1\,000\,300 & 600 & never & 131 & 29 \\
1\,000\,300 & 10\,000 & 3\,000 & 41 & 119 \\
\bottomrule
\end{tabular}
\caption{The trigger on two seconds of the simulator's feed (instrument~1; sell-side add orders at or below the threshold), with a bucket
of 4 tokens refilled every 64 cycles. Verilator, Icarus, the C++20 cycle model and the Python model give the same orders on the same
cycles; orders and rejects add up to the qualifying messages the Python reference finds. Data: \texttt{firm.hwtrade}'s
fixture.}\label{tab:nw:programmable-hardware-ii-trading-designs:runs}
\end{table}

The table shows each control at work. A lower threshold fires less often; the bucket rejects most of the extra triggers of the higher
one, because they come in bursts; the size limit rejects two more; and the kill bit, set at cycle 3\,000, stops every later order. In
the wire-to-wire path of chapter~5 (\cref{fig:nw:programmable-hardware-ii-trading-designs:race}), the hardware path is the cage's fibre,
the card's transceivers (\qty{120}{\nano\second} each way, a model value) and three cycles: \qty{0.82}{\micro\second} at every quantile.
The software path has a median of \qty{4.4}{\micro\second} and a 99th percentile of \qty{35}{\micro\second}. Against a competitor whose
latency has a median of \qty{1}{\micro\second} (a lognormal of spread 0.3, a model), the hardware path wins 75\% of races and the software
path one in ten thousand.

\begin{omfigure}
\begin{tikzpicture}
\begin{groupplot}[group style={group size=2 by 1,horizontal sep=1.8cm},width=6.6cm,height=5.0cm,tick label style={font=\small},
  label style={font=\small},grid=major,grid style={gray!25},table/col sep=comma]
\nextgroupplot[xmin=0,xmax=1,ymode=log,ymin=0.5,ymax=200,xlabel={quantile},ylabel={wire to wire (\unit{\micro\second}, log)},
  ytick={1,10,100},yticklabels={1,10,100},legend columns=2,legend style={font=\footnotesize,at={(1.15,-0.32)},anchor=north,
  draw=none,fill=none,/tikz/every even column/.append style={column sep=8pt}}]
\addplot[omThm,very thick,mark=*,mark size=1.3pt] table[x=q,y=software_us]{figdata/networks/07-programmable-hardware-ii-trading-designs/paths.csv};
\addplot[omDef,very thick,mark=square*,mark size=1.3pt] table[x=q,y=hardware_us]{figdata/networks/07-programmable-hardware-ii-trading-designs/paths.csv};
\legend{software path (chapter~5),hardware path}
\nextgroupplot[xmode=log,xmin=250,xmax=12000,ymin=0,ymax=1.05,xlabel={competitor's median (ns, log)},ylabel={probability of winning},
  xtick={300,1000,3000,10000},xticklabels={300,1\,000,3\,000,10\,000}]
\addplot[omThm,very thick,mark=*,mark size=1.3pt] table[x=median_ns,y=software]{figdata/networks/07-programmable-hardware-ii-trading-designs/race.csv};
\addplot[omDef,very thick,mark=square*,mark size=1.3pt] table[x=median_ns,y=hardware]{figdata/networks/07-programmable-hardware-ii-trading-designs/race.csv};
\end{groupplot}
\end{tikzpicture}

\omcaption{Simulation: wire-to-wire latency of the software path of chapter~5 (Book~13's measured stages, the card and the cage) and of the
hardware path (the cage, the card's transceivers at \qty{120}{\nano\second} each way, a model value, and three cycles at 156.25~MHz). Left:
quantiles; the hardware path is a constant. Right: the probability of beating one competitor whose latency is lognormal with the given
median and a spread of 0.3 (a model). Data: \texttt{fig\_hwtrade.py}.}\label{fig:nw:programmable-hardware-ii-trading-designs:race}
\end{omfigure}

\section{Build: the hardware trigger}\label{bld:nw:programmable-hardware-ii-trading-designs:hwtrade}

\begin{build}
\bfield{Purpose} A complete, tested \omterm{def:nw:programmable-hardware-ii-trading-designs:trigger}{hardware trigger}: feed decoding, risk checks and an \omterm{def:nw:programmable-hardware-ii-trading-designs:template}{order template}, with the three models that make
it trustworthy, and a hardware variant of the firm's wire-to-wire path for chapter~29's plan.
\bfield{Interface} \texttt{hdl/hwt\_trigger.sv} (parameters \texttt{LOC}, \texttt{BURST}, \texttt{REFILL}; registers \texttt{thresh},
\texttt{max\_qty}, \texttt{kill}; outputs the order's id, size, price and a reject counter); \omterm{def:nw:programmable-hardware-i-architecture-and-toolchain:tb}{testbenches} \texttt{tb/hwt\_tb.cpp} and
\texttt{tb/hwt\_tb.sv}; C++20 \texttt{firm::hwtrade::Model}; Python \texttt{firm\_hwtrade}: \texttt{beats\_of}, \texttt{write\_stim},
\texttt{CycleModel}, \texttt{triggers}, \texttt{build\_verilator}, \texttt{build\_icarus}, \texttt{run\_verilator}, \texttt{run\_icarus},
\texttt{parse\_output}, \texttt{wirepath\_stage}.
\bfield{Rules} One order per qualifying message at most; nothing leaves while the kill bit is set; sizes above the limit are rejected, not
clipped; the bucket never lends; the models agree with the HDL register for register.
\bfield{Acceptance tests} \texttt{code/firm/hwtrade/}: the Python model reproduces the fixture's four settings, and every qualifying
message is an order or a reject; Verilator and Icarus print the same orders as the Python model; the C++20 model reproduces the fixture;
the size limit at its boundary; nothing after the kill bit; the bucket's rationing; two cycles from price to order.
\bfield{Stretch} Keep the best offer of the instrument from add, execute and delete messages and fire on it; move the three checks into one
cycle and measure the \omterm{def:nw:programmable-hardware-i-architecture-and-toolchain:timing}{critical path} it would need.
\end{build}

\begin{omsources}
\item STAC, ``STAC-T0'' results with an Exegy and AMD FPGA solution (2024).
\item 17 CFR 240.15c3-5 (the market access rule).
\item Nasdaq, \emph{TotalView-ITCH 5.0} specification (Add Order layout); the simulator's \path{code/firm/exchsim/PROTOCOL.md}.
\end{omsources}

\section{Exercises}

\begin{exercise}[$\star$]\label{exo:nw:programmable-hardware-ii-trading-designs:1}
A packet carries one 36-byte add order behind the 20-byte header and the 2-byte length. In which beat is the price's last byte, and how many
nanoseconds after the packet's first beat does the design emit the order at 156.25~MHz?
\end{exercise}

\begin{exercise}[$\star$]\label{exo:nw:programmable-hardware-ii-trading-designs:2}
With a bucket of 4 tokens refilled every 64 cycles at 156.25~MHz, what is the sustained maximum rate of orders?
\end{exercise}

\begin{exercise}[$\star$]\label{exo:nw:programmable-hardware-ii-trading-designs:3}
How many orders can the chapter's design send before a software monitor that reacts in \qty{50}{\micro\second} sets the kill bit?
\end{exercise}

\begin{exercise}[$\star\star$]\label{exo:nw:programmable-hardware-ii-trading-designs:4}
In \cref{tab:nw:programmable-hardware-ii-trading-designs:runs}, why does raising the threshold from 1\,000\,200 to 1\,000\,300 add 25 orders but
22 rejects?
\end{exercise}

\begin{exercise}[$\star\star$]\label{exo:nw:programmable-hardware-ii-trading-designs:5}
Why does the design reject orders above the size limit instead of sending them at the limit?
\end{exercise}

\begin{exercise}[$\star\star$]\label{exo:nw:programmable-hardware-ii-trading-designs:6}
Name three decisions the chapter's strategy leaves in software, and why each would be hard in hardware.
\end{exercise}

\begin{exercise}[$\star\star\star$]\label{exo:nw:programmable-hardware-ii-trading-designs:7}
\emph{Coding.} With \texttt{firm.hwtrade.CycleModel}, set the bucket to 1 token refilled every 10\,000 cycles and count the orders at threshold
1\,000\,300. Then set the kill bit one cycle before the tenth order of the unrestricted run and count again.
\end{exercise}

\begin{exercise}[$\star\star\star$]\label{exo:nw:programmable-hardware-ii-trading-designs:8}
\emph{Find the flaw.} ``Our hardware fires in 13 nanoseconds, so our wire-to-wire latency is 13 nanoseconds.''
\end{exercise}

\section{Problem: The Trigger That Must Not Fire Twice}

\begin{problem}[{Weekend problem --- a hybrid path, a race and a runaway}]\label{pb:nw:programmable-hardware-ii-trading-designs:1}
A firm moves the decision of a latency-sensitive strategy onto the chapter's \omterm{def:nw:programmable-hardware-ii-trading-designs:trigger}{hardware trigger}. Use the chapter's numbers: the software path's
wire-to-wire median \qty{4.39}{\micro\second} and 99th percentile \qty{35.3}{\micro\second}; the hardware path \qty{0.82}{\micro\second}
(\qty{263}{\nano\second} of cage in, \qty{120}{\nano\second} of card in, three cycles at 156.25~MHz, \qty{120}{\nano\second} of card out,
\qty{295}{\nano\second} of cage out); a competitor with a lognormal latency; a bucket of 4 tokens refilled every 64 cycles.

\medskip\noindent\textbf{Part I --- The path.}
\begin{enumerate}
\item Add up the hardware path's stages.
\item How much of it is the decision itself?
\item What is the hardware path's 99th percentile, and why?
\item By how much does the hardware path beat the software path at the median, and at the 99th percentile?
\end{enumerate}

\medskip\noindent\textbf{Part II --- The race.}
\begin{enumerate}[resume]
\item Against a competitor with a median of \qty{1}{\micro\second}, what are the two paths' probabilities of winning (the chapter's
  simulation)?
\item And against one with a median of \qty{5}{\micro\second}?
\item Which part of the hardware path would you shorten next, and what is left once the decision is gone?
\item What does the race model leave out?
\end{enumerate}

\medskip\noindent\textbf{Part III --- The runaway.}
\begin{enumerate}[resume]
\item A software monitor needs \qty{20}{\micro\second} to see a runaway and set the kill bit. How many orders can leave meanwhile?
\item With a monitor on the card itself, reacting in 16 cycles?
\item Why does the kill bit act on stage 3 rather than only on stage 2?
\item What limits would you add to the card, and which should stay in the broker's layer?
\end{enumerate}

\medskip\noindent\textbf{Part IV --- The verdict.}
\begin{enumerate}[resume]
\item State the \emph{named result}: the two paths' latencies, their probabilities of winning against a competitor with a median of
  \qty{1}{\micro\second}, and the orders the bucket lets out before a \qty{20}{\micro\second} kill.
\item Is the \qty{0.82}{\micro\second} path worth building for a strategy that races competitors at \qty{5}{\micro\second}?
\item What does the firm give up by moving the decision into hardware?
\item How often can the threshold change, and who changes it?
\item Why must the software receive copies of every order the card sends?
\item How would you test the design against a recorded day before it trades?
\item What would a failure of the \omterm{def:nw:programmable-hardware-ii-trading-designs:regmap}{register map}'s write path do, and how would the design detect it?
\item In one sentence: what does hardware decide, and what does software decide?
\end{enumerate}
\end{problem}

\section{Interview questions}

\begin{interviewq}[$\star$ \iqroles{developer}]\label{iq:nw:programmable-hardware-ii-trading-designs:1}
How does an FPGA parse a market-data feed while the packet is still arriving?
\end{interviewq}

\begin{interviewq}[$\star\star$ \iqroles{developer}]\label{iq:nw:programmable-hardware-ii-trading-designs:2}
What risk checks would you put in hardware, and how do you prove they cannot be bypassed?
\end{interviewq}

\begin{interviewq}[$\star\star$ \iqroles{developer, trader}]\label{iq:nw:programmable-hardware-ii-trading-designs:3}
What stays in software in a \omterm{def:nw:programmable-hardware-ii-trading-designs:hybrid}{hybrid trading system}, and why?
\end{interviewq}

\begin{interviewq}[$\star\star$ \iqroles{developer}]\label{iq:nw:programmable-hardware-ii-trading-designs:4}
How do you send an order over TCP from an FPGA?
\end{interviewq}

\begin{interviewq}[$\star\star\star$ \iqroles{developer}]\label{iq:nw:programmable-hardware-ii-trading-designs:5}
How would you verify that the \omterm{def:nw:programmable-hardware-ii-trading-designs:trigger}{hardware trigger} and its software model agree, and what would you do when they disagree?
\end{interviewq}

\begin{interviewq}[$\star\star$ \iqroles{developer, trader}]\label{iq:nw:programmable-hardware-ii-trading-designs:6}
A vendor quotes 14 nanoseconds tick to trade. What questions do you ask?
\end{interviewq}
